\end{lemma}

\begin{proof}
By Lemma {\ref{lem3.2}}, for any $p\in(2, 2^*_s)$, we have
 \begin{equation}\label{5.4}
\|u\|^p_p\leq C{(p,s)}\|(-\Delta)^{s/2}u\|^{\frac{3}{s}(\frac{p}{2}-1)}_2
\|u\|^{\frac{3}{s}(1-\frac{3-2s}{6}p)}_2, \quad \forall u\in H^s(\mathbb{R}^3).
\end{equation}
By the Sobolev inequality \eqref{2.111} and \eqref{5.4}, for $u, v\in  S_{r}(c)$, 
one has
\begin{align*} %\label{5.5}
I_\mu(v)-I_\mu(u)
&= \frac{a}{2}\int_{\mathbb{R}^3}|(-\Delta)^{s/2}v|^2dx
 -\frac{a}{2}\int_{\mathbb{R}^3}|(-\Delta)^{s/2}u|^2dx\\
&\quad +\frac{b}{4}\Big(\int_{\mathbb{R}^3}|(-\Delta)^{s/2}v|^2dx\Big)^2
 -\frac{b}{4}\Big(\int_{\mathbb{R}^3}|(-\Delta)^{s/2}u|^2dx\Big)^2\\
&\quad +\frac{1}{4}\int_{\mathbb{R}^3}\phi^s_v v^2dx-\frac{1}{4}
 \int_{\mathbb{R}^3}\phi^s_u u^2dx+\frac{\mu}{p}\int_{\mathbb{R}^3}|u|^pdx\\
&\quad -\frac{\mu}{p}\int_{\mathbb{R}^3}|v|^pdx+\frac{1}{2^*_s}\int_{\mathbb{R}^3}|u|^{2^*_s}dx-\frac{1}{2^*_s}\int_{\mathbb{R}^3}|v|^{2^*_s}dx\\
&\geq   \frac{a}{2}\Big(\int_{\mathbb{R}^3}|(-\Delta)^{s/2}v|^2dx
 -\int_{\mathbb{R}^3}|(-\Delta)^{s/2}u|^2dx\Big)\\
&\quad +\frac{b}{4}\Big[\Big(\int_{\mathbb{R}^3}|(-\Delta)^{s/2}v|^2dx\Big)^2
 -\Big(\int_{\mathbb{R}^3}|(-\Delta)^{s/2}u|^2dx\Big)^2\Big]\\
&\quad -\frac{1}{4}\Gamma_sC(p_s, s)^{\frac{3+2s}{3}}\|(-\Delta)^{s/2}u\|^{\frac{3-2s}
{s}}_2c^{\frac{6s-3}{s}}\\
&\quad -\frac{\mu}{p}C(p,s)\|(-\Delta)^{s/2}v\|^{p {\delta_{p,s}}}_2 
 c^{\frac{6-p(3-2s)}{2s}}
 -\frac{S^{-\frac{2^*_s}{2}}}{2^*_s}\|(-\Delta)^{s/2}v\|^{2^*_s}_2.
\end{align*}
Let $\|(-\Delta)^{s/2}u\|^2_2\leq K_c$ and $\|(-\Delta)^{s/2}v\|^2_2=2K_c$, 
here $K_c$ will be determined later. Set
$$
\tilde{c}=\Big(\frac{K_c^{\frac{4s-3}{2s}}}{4\Gamma_s C(p_s, s)^{\frac{3+2s}{3}}}a
 \Big)^{\frac{s}{6s-3}},
$$
by a direct computation, we obtain
\begin{equation}\label{5.6}
\begin{split}
&I_\mu(v)-I_\mu(u) \\
&\geq \frac{a}{2}K_c+\frac{3b}{4}K^2_c-\frac{1}{4}\Gamma_s
  C(p_s, s)^{\frac{3+2s}{3}} K^{\frac{3-2s}{2s}}_c 
  \Big(\frac{K_c^{\frac{4s-3}{2s}}}{4\Gamma_s C(p_s, s)^{\frac{3+2s}{3}}}a
   \Big)^{{\frac{s}{6s-3}}\cdot\frac{6s-3}{s}}\\
&\quad -\frac{\mu}{p}C(p,s)(2K_c)^{\frac{p {\delta_{p,s}}}{2}} 
 \Big(\frac{K_c^{\frac{4s-3}{2s}}}{4\Gamma_s C(p_s, s)^{\frac{3+2s}{3}}}a
 \Big)^{\frac{s}{6s-3}\cdot{\frac{6-p(3-2s)}{2s}}}
 -\frac{S^{-\frac{2^*_s}{2}}}{2^*_s}(2K_c)^{\frac{2^*_s}{2}}\\
&\geq  \frac{a}{2}K_c-\frac{a}{16}K_c-\frac{\mu}{p}2^{\frac{p {\delta_{p,s}}}{2}}
 C(p,s)\Big(\frac{1}{4\Gamma_s C(p_s, s)^{\frac{3+2s}{3}}}a
 \Big)^{{\frac{6-p(3-2s)}{2(6s-3)}}}K^{\frac{(4s-3)(6-p(3-2s))}{4s(6s-3)}}_c 
  K^{\frac{3(p-2)}{4s}}_c\\
&\quad -\frac{S^{-\frac{2^*_s}{2}}}{2^*_s}2^{\frac{2^*_s}{2}}(K_c)^{\frac{2^*_s}{2}}\\
&= \frac{7}{16}aK_c-\frac{\mu 2^{\frac{3(p-2)}{4s}}C(p,s)a^{{\frac{6-p(3-2s)}{2(6s-3)}}}}{p(4\Gamma_s C(p_s, s)^{\frac{3+2s}{3}})^{\frac{6-p(3-2s)}{2(6s-3)}}}K^{\gamma_1}_c K_c-\frac{2^{\frac{2^*_s}{2}}}{2^*_sS^{\frac{2^*_s}{2}}}K^{\frac{2^*_s-2}{2}}_c K_c\\
&\geq   \frac{5}{16}aK_c>0,
\end{split}
\end{equation}
where 
\[
\gamma_1=\frac{[6-p(3-2s)][4s-3]+[3(p-2)-4s][6s-3]}{4s(6s-3)}.
\] 
If we take
$$
K_c=\min\Big\{ \Big(\frac{{p(4\Gamma_s C(p_s, s)^{\frac{3+2s}{3}}) }
 ^{\frac{6-p(3-2s)}{2(6s-3)}}}{16 \mu 2^{\frac{3(p-2)}{4s}}C(p,s)
 a^\frac{6-p(3-2s)}{2(6s-3)}} a\Big)^{\gamma_2}, 
 \Big(\frac{2^*_s S^{\frac{2^*_s}{2}}}{2^{\frac{2^*_s}{2}}16} a
  \Big)^{\frac{2}{2^*_s-2}}\Big\},
$$
with 
\[
\gamma_2=\frac{4s(6s-3)}{[6-p(3-2s)][4s-3]+[3(p-2)-4s][6s-3]},
\]
then, by \eqref{5.6} we deduce that \eqref{5.3} holds. The proof is complete.
\end{proof}

From Lemma \ref{lem5.2}  we can deduce the following conclusion.

\begin{corollary}\label{cor5.1}  
Let $K_c$ and $\tilde{c}$ be given in Lemma \ref{lem5.2}, and $u\in  S_{r}(c)$ with 
$\|(-\Delta)^{s/2}u\|^2_2 \leq K_c$, then $I_\mu(u)>0$.
\end{corollary}

\begin{proof}
A direct computation shows that
\begin{align*} %\label{5.7}
I_\mu(u)
&\geq  \frac{a}{2}\|(-\Delta)^{s/2}u\|^2_2+\frac{b}{4}\|(-\Delta)^{s/2}u\|^4_2
 -\frac{\mu}{p}C(p,s)
c^{\frac{6-p(3-2s)}{2s}}\|(-\Delta)^{s/2}u\|^{\frac{3(p-2)}{2s}}_2\\
&\quad -\frac{S^{-\frac{2^*_s}{2}}}{2^*_s}\|(-\Delta)^{s/2}u\|^{2^*_s}_2>0,
\end{align*}
if $\|(-\Delta)^{s/2}u\|^2_2 \leq K_c$, and the conclusion holds.
\end{proof}

Next, we study the characterizations of the mountain pass levels for 
$I(u,\theta)$ and $I_\mu(u)$. We denote the closed set
$I^d_\mu:=\{u\in  S_{r}(c):I_\mu(u)\leq d\}$.

\begin{proposition}\label{pro5.1} 
Assuming that  $\frac{8s}{3}+2<p< 2^*_s$,  we define
$$
\tilde{c}_\mu (c):=\inf_{\tilde{\gamma}\in \tilde{\Gamma}}\max_{t\in[0,1]}
I(\tilde{\gamma}(t)),~~~{c}_\mu (c):=\inf_{{\gamma}\in{\Gamma}}
\max_{t\in[0,1]}I_{\mu}({\gamma}(t)),
$$
where
\begin{gather*}
\tilde{\Gamma}_c=\{\tilde{\gamma}\in C([0,1],  S_{r}(c)\times \mathbb{R}): 
\tilde{\gamma}(0)\in (\mathcal{A}_c, 0), \tilde{\gamma}(1)\in (I^0_\mu, 0)\},\\
{\Gamma}_c=\{{\gamma}\in C([0,1],  S_{r}(c)): {\gamma}(0)\in \mathcal{A}_c,
 {\gamma}(1)\in I^0_\mu\}.
\end{gather*}
Then we have 
$\tilde{c}_\mu (c)={c}_\mu (c)>0$.
\end{proposition}

\begin{proof}
 On the one hand, for any $\tilde{\gamma}\in \tilde{\Gamma}_c$, we can write  it into 
$$
\tilde{\gamma}(t)=(\tilde{\gamma}_1(t), \tilde{\gamma}_2(t))\in S_r(c)\times \mathbb{R}.
$$
We set ${\gamma}(t)=\tilde{\gamma}_2(t)\star \tilde{\gamma}_1(t)$, then 
${\gamma}\in {\Gamma}_c$, and
$$
 \max_{t\in[0,1]} I(\tilde{\gamma}(t))=\max_{t\in[0,1]} I_\mu(\tilde{\gamma}_2(t)
 \star \tilde{\gamma}_1(t))=\max_{t\in[0,1] } I_\mu({\gamma}(t)),
$$
which implies $\tilde{c}_\mu (c)\geq {c}_\mu (c)>0$, using Corollary \ref{cor5.1}. 
On the other hand, for any ${\gamma}\in {\Gamma}_c$, if we set
$\tilde{\gamma}(t)= ({\gamma}(t), 0)$, then we obtain 
$\tilde{\gamma}\in \tilde{\Gamma}_c$ and
$$
\max_{t\in[0,1]} I(\tilde{\gamma}(t))
=\max_{t\in[0,1] } I_\mu({\gamma}(t)).
$$
This infers that  $\tilde{c}_\mu (c)\leq {c}_\mu (c)$. So,  
$\tilde{c}_\mu (c)= {c}_\mu (c)>0$.
\end{proof}

Next, we show the existence of the $(PS)_{{c}_\mu (c)}$-sequence for $I(u, \theta)$ 
on $ S_{r}(c)\times \mathbb{R}\subset \mathbb{H}$. 
It is obtained by a standard argument using Ekeland's variational principle and
constructing pseudo-gradient flow, see \cite[Proposition 2.2]{J7888}.

\begin{lemma}\label{lem5.4} 
Let $\{\tilde{h}_n\}\subset \tilde{\Gamma}_c $ satisfy that
$$
\max_{t\in[0,1]}I(\tilde{h}_n(t))\leq \tilde{c}_\mu (c)+\frac{1}{n},
$$
then there exists a sequence $\{(v_n, \theta_n)\}\subset S_{r}(c)\times \mathbb{R}$ 
such that
\begin{itemize}
\item[(i)] $I{(v_n, \theta_n)}\in [\tilde{c}_\mu (c)-\frac{1}{n}, 
 \tilde{c}_\mu (c)+\frac{1}{n}]$,
 
\item[(ii)] $\min_{t\in[0,1]}\|(v_n, \theta_n)-\tilde{h}_n(t)\|_{\mathbb{H}}
\leq \frac{1}{\sqrt{n}}$; and

\item[(iii)] $\|(I|_{S_{r}(c)\times \mathbb{R}})'(v_n, \theta_n)\|
 \leq \frac{2}{\sqrt{n}}$, that is,
$$
|\langle I'(v_n, \theta_n), z\rangle|_{\mathbb{H}^{-1}\times \mathbb{H}}
\leq \frac{2}{\sqrt{n}} \|z\|_{\mathbb{H}},
$$
for all
$$ 
{z\in {\tilde{T}_{(v_n, \theta_n)}}}=\{(z_1, z_2)\in \mathbb{H}: 
{\langle v_n, z_1\rangle}_{L^2}=0\}.
$$
\end{itemize}
\end{lemma}


It follows from the above proposition, we can obtain a special 
$(PS)_{{c}_\mu (c)}$-sequence for $I_\mu(u)$ on $S_{r}(c)\subset H^s_r(\mathbb{R}^3)$.

\begin{lemma}\label{lem5.5} 
Under the  assumption $2+\frac{8s}{3}<p<2^*_s$, there exists a sequence
$\{u_n\} \subset S_{r}(c)$ such that
\begin{itemize}
\item[(1)] $I_\mu(u_n)\to c_\mu(c) ~as~ n\to \infty$;

\item[(2)] $P_\mu(u_n)\to 0 ~as~ n\to \infty$;

\item[(3)] $(I_\mu|_{S_{r}(c)})'(u_n)\to 0$  as $n\to \infty$, i.e., 
$\langle I'_\mu(u_n), z\rangle_{{H^{-s}_r} \times H^s_r} \to 0$, 
uniformly for all $z$ satisfying
$$
\|z\|_{H^s_r}\leq 1,\quad \text{where } z\in T_{u_n} 
:=\{z\in {H^s_r(\mathbb{R}^3)}:~\langle u_n,z \rangle_{L^2}=0\}.
$$
\end{itemize}
\end{lemma}

\begin{proof}
Let $\{h_n\}\subset \Gamma_c$ satisfy 
\begin{equation}\label{5.8}
\max_{t\in[0.1]}I_\mu(h_n(t))\leq {c}_\mu(c)+\frac{1}{n},
\end{equation}
we define $\tilde{h}_n(t)=(h_n(t) , 0), \forall t\in[0,1]$. 
It is easy to see that $\tilde{h}_n \in \tilde{\Gamma}_c$ and 
$I_\mu(h_n(t))=I(\tilde{h}_n(t))$. By Proposition \ref{pro5.1}, we have 
$\tilde{c}_\mu(c)={c}_\mu(c)$, then it follows from \eqref{5.8} that
\[  %\label{5.9}
\max_{t\in[0.1]}I(\tilde{h}_n(t))\leq \tilde{c}_\mu(c)+\frac{1}{n}.
\]
It follows from Lemma \ref{lem5.4} that, there exists a sequence
$\{(v_n, \theta_n)\} \subset S_{r}(c)\times \mathbb{R}$ such that as $n\to \infty$,
one has
\begin{gather}\label{5.9}
I(v_n, \theta_n)\to c_\mu(c), ~\theta_n\to 0, \\
\label{5.10}
(I|_{S_{r}(c)\times{\mathbb{R}}})'(v_n, \theta_n)\to 0.
\end{gather}
Set $u_n=\theta_n \star  v_n$. Then, $I_\mu(u_n)=I(v_n, \theta_n)$,
and by \eqref{5.9}, item (1) holds.  To prove conclusion (2), we utilize
\begin{align*}% \label{5.11}
\partial_\theta I(v_n, \theta_n)
&= ase^{2s\theta_n}\|(-\Delta)^{s/2}v_n\|^2_2+bse^{4s\theta_n}\|(-\Delta)^{s/2}v_n\|^4_2+\frac{3-2s}{4}e^{(3-2s)\theta_n}
\int_{\mathbb{R}^3}\phi^s_{v_n}{v^2_n}dx\\
&\quad -\frac{\mu}{p}\frac{3(p-2)}{2}e^{\frac{3(p-2)}{2}\theta n}\int_{\mathbb{R}^3}|v_n|^pdx-\frac{3(2^*_s-2)}{2 2^*_s}e^{\frac{3(2^*_s-2)}{2}\theta n}\int_{\mathbb{R}^3}|v_n|^{2^*_s}dx\\
&= s(a+b\|(-\Delta)^{s/2}u_n\|^2_2)\|(-\Delta)^{s/2}u_n\|^2_2+\frac{3-2s}{4}\int_{\mathbb{R}^3}\phi^s_{u_n}{u^2_n}dx\\
&\quad -\frac{3(p-2)}{2p}\mu \int_{\mathbb{R}^3}|u_n|^pdx-s\int_{\mathbb{R}^3}|u_n|^{2^*_s}dx\\
&= P_\mu(u_n),
\end{align*}
which implies item (2) by \eqref{5.10}. To show item (3), we set $z_n\in T_{u_n}$.
Then,
\begin{align*}%\label{5.12}
\langle I'_\mu(u_n),z_n\rangle
&= a\int_{\mathbb{R}^3}\int_{\mathbb{R}^3}
 \frac{(u_n(x)-u_n(y))(z_n(x)-z_n(y))}{|x-y|^{3+2s}}\,dx\,dy\\
&\quad +b\int_{\mathbb{R}^3}\int_{\mathbb{R}^3}\frac{(u_n(x)
 -u_n(y))^2}{|x-y|^{3+2s}}\,dx\,dy
 \int_{\mathbb{R}^3}\int_{\mathbb{R}^3}
 \frac{(u_n(x)-u_n(y))(z_n(x)-z_n(y))}{|x-y|^{3+2s}}\,dx\,dy  \\
&\quad +\int_{\mathbb{R}^3}\phi^s_{u_n}{u_n}z_ndx
 -\mu\int_{\mathbb{R}^3}|u_n|^{p-2}u_n z_ndx-\int_{\mathbb{R}^3}|u_n|^{2^*_s-2}u_n z_ndx\\
&=  a\int_{\mathbb{R}^3}\int_{\mathbb{R}^3}\frac{(e^{\frac{3\theta_n}{2}}
 v_n(e^{\theta_n} x)-e^{\frac{3\theta_n}{2}}v_n(e^{\theta_n} y))(z_n(x)
 -z_n(y))}{|x-y|^{3+2s}}\,dx\,dy\\
&\quad +b\int_{\mathbb{R}^3}\int_{\mathbb{R}^3}
 \frac{(e^{\frac{3\theta_n}{2}}v_n(e^{\theta_n} x)-e^{\frac{3\theta_n}{2}}
 v_n(e^{\theta_n} y))^2}{|x-y|^{3+2s}}\,dx\,dy \\
 &\quad\times \int_{\mathbb{R}^3}
 \int_{\mathbb{R}^3}\frac{(e^{\frac{3\theta_n}{2}}v_n(e^{\theta_n } x)
 -e^{\frac{3\theta_n}{2}}v_n(e^{\theta_n} y))(z_n(x)-z_n(y))}{|x-y|^{3+2s}}\,dx\,dy\\
&\quad +e^{\frac{3-4s}{2}\theta_n}\int_{\mathbb{R}^3}\phi^s_{v_n}{v_n}
 z_n(e^{-\theta_n}x)dx-\mu e^{\frac{3(p-3)}{2}\theta_n}\int_{\mathbb{R}^3}
 |v_n|^{p-2}v_n z_n(e^{-\theta_n}x) dx\\
&\quad -e^{\frac{3(2^*_s-3)}{2}\theta_n}\int_{\mathbb{R}^3}|v_n|^{2^*_s-2}
 v_n z_n(e^{-\theta_n}x) dx\\
&= e^{2s\theta_n}a\int_{\mathbb{R}^3}\int_{\mathbb{R}^3}\frac{(v_n(x)-v_n(y))
 e^{-\frac{3}{2}\theta_n}(z_n(e^{-\theta_n} x)-z_n(e^{-\theta_n}y))}
  {|x-y|^{3+2s}}\,dx\,dy\\
&\quad +e^{4s\theta_n}b\int_{\mathbb{R}^3}\int_{\mathbb{R}^3}
 \frac{(v_n( x)-v_n( y))^2}{|x-y|^{3+2s}}\,dx\,dy \\
 &\quad\times  \int_{\mathbb{R}^3}\int_{\mathbb{R}^3}\frac{(v_n( x)-v_n( y))
 e^{-\frac{3}{2}\theta_n}(z_n(e^{-\theta_n} x)
 -z_n(e^{-\theta_n}y))}{|x-y|^{3+2s}}\,dx\,dy  \\
&\quad +e^{\frac{3-4s}{2}\theta_n}\int_{\mathbb{R}^3}\phi^s_{v_n}{v_n}z_n
 (e^{-\theta_n}x)dx-\mu e^{\frac{3(p-3)}{2}\theta_n}
 \int_{\mathbb{R}^3}|v_n|^{p-2}v_n z_n(e^{-\theta_n}x) dx\\
&\quad -e^{\frac{3(2^*_s-3)}{2}\theta_n}\int_{\mathbb{R}^3}|v_n|^{2^*_s-2}
v_n z_n(e^{-\theta_n}x) dx.
\end{align*}
Denoting $\tilde{z}_n(x)=e^{-\frac{3\theta_n}{2}}z_n(e^{-\theta_n}x)$,
 we obtain
$$
\langle I'_\mu(u_n),z_n\rangle_{{H^{-s}_r}\times H^s_r}
=\langle I'(v_n, \theta_n), (\tilde{z}_n, 0) \rangle_{\mathbb{H}^{-1}\times \mathbb{H}}.
$$
It is easy to check that
\[  %\label{5.12}
 \langle v_n, \tilde{z}_n\rangle_{L^2}=\int_{\mathbb{R}^3} v_n(x) e^{-\frac{3\theta_n}{2}}z_n(e^{-\theta_n}x)dx
 =\int_{\mathbb{R}^3} v_n(e^{\theta_n}x) e^{\frac{3\theta_n}{2}}z_n(x)dx
 =\int_{\mathbb{R}^3}u_n(x)z_n(x)dx=0.
\]
Therefore, $(\tilde{z}_n, 0)\in \tilde{T}_{(v_n, \theta_n)}$.
On the other hand,
$$
\|(\tilde{z}_n,0)\|^2_{\mathbb{H}}=\|\tilde{z}_n\|^2_{{H^s_r}}=\|z_n\|^2_2
+e^{-2s\theta_n}\|z_n\|^2_{D^{s,2}}\leq C\|z_n\|^2_{H^s_r},
$$
where the last inequality follows by $\theta_n\to 0$. Consequently, we conclude
item (3). The proof is complete.
\end{proof}


\begin{lemma}\label{lem5.6} 
The $(PS)$ sequence $\{u_n\} \subset S_{r}(c)$ for $I_\mu(u)$ with the level 
$c_\mu(c)$ mentioned in Lemma \ref{lem5.5} is bounded in $H^s_r(\mathbb{R}^3)$.
\end{lemma}

\begin{proof}  
From Lemma  \ref{lem5.5} (1), we see that $I_\mu(u)$ is bounded. 
In fact, by $P_\mu(u_n)\to 0$ as $n\to \infty$, one obtains
$$
|(1+2s)I_\mu(u_n)+P_\mu(u_n)|\leq 3c_\mu(c),
$$
which implies that
\begin{equation}\label{5.13}
\begin{split}
-3c_\mu(c)&\leq \frac{1+4s}{2}a\|(-\Delta)^{s/2}u_n\|^2_2+\frac{1+6s}{4}b
\|(-\Delta)^{s/2}u_n\|^4_2+\int_{\mathbb{R}^3}\phi^s_{u_n}u^2_ndx\\
&\quad -\mu(\frac{1+2s}{p}+s\delta_{p,s})\int_{\mathbb{R}^3}|u_n|^pdx
 -(\frac{1+2s}{2^*_s}+s)\int_{\mathbb{R}^3}|u_n|^{2^*_s}dx.
\end{split}
\end{equation}
In view of the boundedness of $I_\mu(u)$, we have
\begin{equation}\label{5.14}
\begin{split}
&a\|(-\Delta)^{s/2}u_n\|^2_2+\frac{b}{2}\|(-\Delta)^{s/2}u_n\|^4_2
 +\frac{1}{2}\int_{\mathbb{R}^3}\phi^s_{u_n}u^2_ndx\\
&\leq 6c_\mu(c)+\frac{2\mu}{p}\int_{\mathbb{R}^3}|u_n|^pdx
 +\frac{2}{2^*_s}\int_{\mathbb{R}^3}|u_n|^{2^*_s}dx.
\end{split}
\end{equation}
By \eqref{5.13} and \eqref{5.14}, we obtain
\begin{align*} %\label{5.15}
\mu \frac{(p\delta_{p,s}-4)s}{p}\int_{\mathbb{R}^3}|u_n|^pdx
+\frac{(2^*_s-4)s}{2^*_s}\int_{\mathbb{R}^3}|u_n|^{2^*_s}dx
+\frac{(6s-3)}{4}\int_{\mathbb{R}^3}\phi^s_{u_n}u^2_ndx\leq 3(2+6s)c_\mu(c).
\end{align*}
Note that $s\in(\frac{3}{4}, 1)$, $p>\frac{8s}{3}+2$, we have that 
$p\delta_{p,s}-4>0$, and so
$$
\int_{\mathbb{R}^3}|u_n|^pdx, \int_{\mathbb{R}^3}|u_n|^{2^*_s}dx 
\quad\text{and}\quad
\int_{\mathbb{R}^3}\phi^s_{u_n}u^2_ndx 
$$
are all bounded. Thus, $\|(-\Delta)^{s/2}u_n\|_2\leq R_3$ for some $R_3>0$ 
independently on $n\in \mathbb{N}$. Since $\{u_n\}\subset S_{r}(c)$, we see that 
$\{u_n\}$ is bounded in $H^s_r(\mathbb{R}^3)$.  
This completes the proof.
\end{proof}

Now, we set the functional $\Phi:H^s_r(\mathbb{R}^3)\to \mathbb{R}$ as
$$
\Phi(u)=\frac{1}{2}\int_{\mathbb{R}^3} |u|^2dx,
$$
then $S_{r}(c)=\Phi^{-1}(\frac{c^2}{2})$. As a result, it can be derived from 
\cite[Proposition 5.12]{W6888} that there is a sequence 
$\{\lambda_n\}\subset \mathbb{R}$ such that
$$
I'_\mu(u_n)-\lambda_n \Phi'(u_n)\to 0,~in ~H^{-s}_r(\mathbb{R}^3) ~as~n\to \infty.
$$
That is, in $H^{-s}_r(\mathbb{R}^3)$, we have
\begin{equation}\label{5.16}
\left(a+b\|(-\Delta)^{s/2}u_n\|^2_2\right)(-\Delta)^s u_n+\phi^s_{u_n}u_n
-\mu |u_n|^{p-2}u_n-|u_n|^{2^*_s-2}u_n
=\lambda_n u_n+o_n(1).
\end{equation}
Therefore, for any $\varphi \in H^s_r(\mathbb{R}^3)$, one has
\begin{equation}\label{5.166}
\begin{split}
&\left(a+b\|(-\Delta)^{s/2}u_n\|^2_2\right)
\int_{\mathbb{R}^3}(-\Delta)^{s/2} u_n(-\Delta)^{s/2} \varphi dx
+\int_{\mathbb{R}^3}\phi^s_{u_n}u_n\varphi dx \\
&-\int_{\mathbb{R}^3}\mu |u_n|^{p-2}u_n\varphi dx
 -\int_{\mathbb{R}^3}|u_n|^{2^*_s-2}u_n\varphi dx \\
&=\lambda_n \int_{\mathbb{R}^3}u_n\varphi dx +o_n(1).
\end{split}
\end{equation}

Next, we study the asymptotical behavior of the mountain pass level value 
$c_\mu(c)$ as $\mu \to +\infty$, and the properties of  the $(PS)_{c_\mu(c)}$-sequence 
$\{u_n\}\subset S_{r}(c)$ as $n\to +\infty$.

\begin{lemma}\label{lem5.7} 
The limit $\lim_{\mu\to +\infty}c_\mu(c)=0 $ holds.
\end{lemma}

\begin{proof} 
Recall Lemma {\ref{lem5.1}} and Corollary {\ref{cor5.1}}, we see that for fixed 
$u_0\in S_{r}(c)$, there exists two constants
$\theta_1, \theta_2$ satisfying $\theta_1<0< \theta_2$ such that 
$u_1=\theta_1 \star u_0 \in \mathcal{A}_c$ and 
$I_\mu(u_2)=I_\mu((\theta_2 \star u_0 ))<0$. Then, we can define a path
$$
\eta_0: t \in [0,1]\to ((1-t)\theta_1+t \theta_2)\star u_0 \in \Gamma_c.
$$
Therefore,
\begin{align*}
0<c_\mu(c)
&\leq  \max_{t\in[0,1]} I_\mu(\eta_0(t))\\
&\leq  \max_{r\geq 0} \Big\{\frac{a}{2}r^{2s}\|(-\Delta)^{s/2}u_0\|^2_2+\frac{b}{4}r^{4s}\|(-\Delta)^{s/2}u_0\|^4_2\\
&\quad +\frac{1}{4}r^{3-2s}\int_{\mathbb{R}^3}\phi^s_{u_0}u^2_0dx
 -\frac{\mu}{p}r^{\frac{3(p-2)}{2}}\int_{\mathbb{R}^3}|u_0|^pdx\Big\}\\
&:=\max_{r\geq 0}g(r).
\end{align*}
Note that $\frac{3(p-2)}{2}>4s>2s>3-2s$, we have that $\lim_{r\to 0^+}g(r)=0^+$, 
$\lim_{r\to +\infty}g(r)=-\infty$. Then, there exists a unique maximum point $r_0>0$ 
such that $\max_{r\geq 0}g(r)=g(r_0)>0$. Hence, we distinguish two cases: $r_0\geq1$ 
and $0 \leq r_0<1$.

If $r_0\geq1$, then by $s\in(\frac{3}{4}, 1)$, we have
\begin{align*}
\max_{t\in[0,1]} I_\mu(\eta_0(t))&\leq  g(r_0)\\
&\leq  \Big\{\frac{a}{2}r^{4s}_0\|(-\Delta)^{s/2}u_0\|^2_2
 +\frac{b}{4}r^{4s}_0\|(-\Delta)^{s/2}u_0\|^4_2\\
&\quad +\frac{1}{4}r^{4s}_0\int_{\mathbb{R}^3}\phi^s_{u_0}u^2_0dx
 -\frac{\mu}{p}r^{\frac{3(p-2)}{2}}_0\int_{\mathbb{R}^3}|u_0|^pdx\Big\}\\
&\leq  \max_{r\geq 0}\Big\{3\max\Big\{\frac{a}{2}\|(-\Delta)^{s/2}u_0\|^2_2, 
 \frac{b}{4}\|(-\Delta)^{s/2}u_0\|^4_2,
 \frac{1}{4}\int_{\mathbb{R}^3}\phi^s_{u_0}u^2_0dx  \Big\}r^{4s}\\
&\quad -\frac{\mu}{p}r^{\frac{3(p-2)}{2}}\int_{\mathbb{R}^3}|u_0|^pdx\Big\}\\
&= 3m(r_{\rm max})^{4s}-\frac{\mu}{p}n(r_{\rm max})^{\frac{3(p-2)}{2}}\\
&= \frac{m(3p-6-8s)}{p-2}\left[\frac{8psm}{(p-2)\mu n}\right]^{\frac{8s}{3p-6-8s}},
\end{align*}
where $r_{\rm max}=\left[\frac{8psm}{(p-2)\mu n}\right]^{\frac{2}{3p-6-8s}}, 
~m=\max\left\{\frac{a}{2}\|(-\Delta)^{s/2}u_0\|^2_2, \frac{b}{4}\|(-\Delta)^{s/2}u_0\|^4_2,
\frac{1}{4}\int_{\mathbb{R}^3}\phi^s_{u_0}u^2_0dx  \right\},$ $n=\int_{\mathbb{R}^3}|u_0|^pdx.$
Therefore, for $\frac{8s}{3}+2<p<2^*_s$, we have a positive constant $\tilde{C}$ independent of $\mu$ such that
$$c_\mu (c)\leq \tilde{C}{\mu}^{-\frac{8s}{3p-6-8s}}\to 0, ~~as~\mu\to +\infty.$$


If $0 \leq r_0<1$, we infer to
\begin{align*}
\max_{t\in[0,1]} I_\mu(\eta_0(t))&\leq  g(r_0)\\
&\leq  \Big\{\frac{a}{2}r^{3-2s}_0\|(-\Delta)^{s/2}u_0\|^2_2
 +\frac{b}{4}r^{3-2s}_0\|(-\Delta)^{s/2}u_0\|^4_2\\
&\quad +\frac{1}{4}r^{3-2s}_0\int_{\mathbb{R}^3}\phi^s_{u_0}u^2_0dx
 -\frac{\mu}{p}r^{\frac{3(p-2)}{2}}_0\int_{\mathbb{R}^3}|u_0|^pdx\Big\}\\
&\leq  \max_{r\geq 0}\Big\{3\max\Big\{\frac{a}{2}\|(-\Delta)^{s/2}u_0\|^2_2, 
 \frac{b}{4}\|(-\Delta)^{s/2}u_0\|^4_2,
\frac{1}{4}\int_{\mathbb{R}^3}\phi^s_{u_0}u^2_0dx  \Big\}r^{3-2s}\\
&\quad -\frac{\mu}{p}r^{\frac{3(p-2)}{2}}\int_{\mathbb{R}^3}|u_0|^pdx\Big\}\\
&= 3m(r_{\rm max})^{3-2s}-\frac{\mu}{p}n(r_{\rm max})^{\frac{3(p-2)}{2}}\\
&= \frac{m(3p+4s-12)}{p-2}\Big[\frac{2pm(3-2s)}{(p-2)\mu n}\Big]^{\frac{6-4s}{3p+4s-12}},
\end{align*}
where $r_{\rm max}=[\frac{2pm(3-2s)}{(p-2)\mu n}]^{\frac{2}{3p+4s-12}}$.
Therefore, for $2+\frac{8s}{3}<p<2^*_s$, and $s\in{(\frac{3}{4}, 1)}$, we can deduce 
that $3p+4s-12>0$, then there exists a positive constant $\bar{C}$ independent of
 $\mu$ such that
$$
c_\mu (c)\leq \bar{C}{\mu}^{-\frac{6-4s}{3p+4s-12}}\to 0, ~~as~\mu\to +\infty.
$$
This completes the proof.
\end{proof}

\begin{lemma}\label{lem5.8} 
There exists  a constant $C=C(p,s)>0$ such that
\begin{gather*}
\limsup_{n\to \infty} \int_{\mathbb{R}^3}F(u_n)dx\leq C c_\mu(c), \quad
\limsup_{n\to \infty} \int_{\mathbb{R}^3}f(u_n){u_n}dx\leq C c_\mu(c), \\
\limsup_{n\to \infty} \int_{\mathbb{R}^3}\phi^s_{u_n}u^2_ndx\leq C c_\mu(c),\quad 
\limsup_{n\to \infty} \int_{\mathbb{R}^3}|(-\Delta)^{s/2}u_n|^2dx\leq C c_\mu(c),\\
\limsup_{n\to \infty} (\int_{\mathbb{R}^3}|(-\Delta)^{s/2}u_n|^2dx)^2\leq C c_\mu(c),
\end{gather*}
where $f(u)=\mu|u|^{p-2}u+|u|^{2^*_s-2}u$, $F(u)=\int^u_0 f(s)ds$.
\end{lemma}

\begin{proof} 
Since $I_\mu(u_n)\to c_\mu(c)$ and $P_\mu(u_n)\to 0$ as $n \to \infty$, one obtains
\begin{equation}\label{5.17}
\begin{split}
&3c_\mu(c)+o_n {(1)}=3I_\mu(u_n)+P_\mu(u_n)\\
&= \frac{3+2s}{2}a\int_{\mathbb{R}^3}|(-\Delta)^{s/2}u_n|^2dx
+\frac{3+4s}{4}b\Big(\int_{\mathbb{R}^3}|(-\Delta)^{s/2}u_n|^2dx\Big)^2\\
&\quad +\frac{3-s}{2}\int_{\mathbb{R}^3}\phi^s_{u_n} u^2_ndx
 -\frac{3}{2}\int_{\mathbb{R}^3}f(u_n)u_ndx\\
&\leq  \frac{3+4s}{2}a\int_{\mathbb{R}^3}|(-\Delta)^{s/2}u|^2dx
+\frac{3+4s}{4}b\Big(\int_{\mathbb{R}^3}|(-\Delta)^{s/2}u|^2dx\Big)^2\\
&\quad +\frac{3-s}{2}\int_{\mathbb{R}^3}\phi^s_{u_n} u^2_ndx-\frac{3}{2}\int_{\mathbb{R}^3}f(u_n)u_ndx\\
&= (3+4s)\Big(-\frac{1}{4}\int_{\mathbb{R}^3}\phi^s_{u_n} u^2_ndx
 +\int_{\mathbb{R}^3}F(u_n)dx+c_\mu(c)+o_n(1)\Big)\\
&\quad +\frac{3-s}{2}\int_{\mathbb{R}^3}\phi^s_{u_n} u^2_ndx
 -\frac{3}{2}\int_{\mathbb{R}^3}f(u_n)u_ndx\\
&= (3+4s)\Big(\int_{\mathbb{R}^3}F(u_n)dx+c_\mu(c)+o_n(1)\Big)\\
&\quad -\frac{6s-3}{4}\int_{\mathbb{R}^3}\phi^s_{u_n} u^2_ndx-\frac{3}{2}\int_{\mathbb{R}^3}f(u_n)u_ndx\\
&\leq  (3+4s)\Big(\int_{\mathbb{R}^3}F(u_n)dx+c_\mu(c)+o_n(1)\Big)
 -\frac{3}{2}\int_{\mathbb{R}^3}f(u_n)u_ndx\\
&\leq  (3+4s)(c_\mu(c)+o_n(1))-\frac{3}{2}p\int_{\mathbb{R}^3}F(u_n)dx
 +(3+4s)\int_{\mathbb{R}^3}F(u_n)dx.
\end{split}
\end{equation}
Hence,
$$
4sc_\mu(c)+o_n(1)\geq \frac{3p-6-8s}{2}\int_{\mathbb{R}^3}F(u_n)dx,
$$
which implies that
\begin{equation}\label{5.18}
\limsup_{n\to \infty} \int_{\mathbb{R}^3}F(u_n)dx
\leq \frac{8s}{3p-6-8s}c_\mu(c) 
\leq C c_\mu(c)
\end{equation}
and then
\begin{equation}\label{5.19}
\limsup_{n\to \infty} \int_{\mathbb{R}^3}f(u_n)u_ndx \leq C c_\mu(c).
\end{equation}
Then, from \eqref{5.17}-\eqref{5.19}, one has
\begin{align*} %\label{5.20}
&\limsup_{n\to \infty}\Big\{\frac{3+2s}{2}a
 \int_{\mathbb{R}^3}|(-\Delta)^{s/2}u_n|^2dx
 +\frac{3+4s}{4}b(\int_{\mathbb{R}^3}|(-\Delta)^{s/2}u_n|^2dx)^2
 +\frac{3-s}{2}\int_{\mathbb{R}^3}\phi^s_{u_n} u^2_ndx\Big\}\\
&= \limsup_{n\to \infty}\Big\{\frac{3}{2}\int_{\mathbb{R}^3}f(u_n)u_ndx+3c_\mu(c)\Big\}\\
&\leq  C c_\mu(c).
\end{align*}
Consequently, the proof is complete.
\end{proof}

\begin{lemma}\label{rem5.91} 
Let $\{u_n\} \subset S_r(c)$ be the  $(PS)$ sequence for the constrained functional
 $I_\mu|_{S_r(c)}$ at level 
 $c_\mu(c)\in \left(0, \Big(\frac{1}{4}-\frac{1}{2^*_s}\Big)(aS)^{\frac{3}{2s}}\right)$ 
with $P_\mu(u_n)\to 0$ as $n\to \infty$. Then
\begin{itemize}
\item[(i)] $\{\lambda_n\}$ is bounded in $\mathbb{R}$, and 
$\limsup_{n\to \infty}|\lambda_n|\leq \frac{C}{c^2}c_\mu(c)$ has the  
estimation
$$
\lambda_n=\frac{1}{c^2}\Big[\frac{6s-3}{4s}\int_{\mathbb{R}^3}\phi^s_{u_n} u^2_ndx
+\mu \frac{p(3-2s)-6}{2ps}\int_{\mathbb{R}^3}|u_n|^pdx\Big]+o_n(1).
$$
Moreover, there exists some $\mu^*_2:=\mu^*_2(c)>0$ such that 
$\lim_{n\to +\infty}\lambda_n=\lambda<0$, if $\mu> \mu^*_2$ large;

\item[(ii)] there exist $u\in S_r(c)$ such that, up to a subsequence, 
$u_n\to u$ strongly in $H^s_r(\mathbb{R}^3)$ as  $\mu> \mu^*_2$ large and 
$u$ is a solution of  system \eqref{lab1.1} for some $\lambda<0$.
\end{itemize}
\end{lemma}

\begin{proof}  We split the proof into three steps.
\smallskip

\noindent\textbf{Step 1.} 
We assert that  $u\not\equiv 0$.
From Lemma {\ref{lem5.6}}, we  know that $\{u_n\}$ is a  bounded $(PS)$ sequence 
for $I_\mu$ in $H^s_r(\mathbb{R}^3)$, and by Lemma {\ref{lem2.1311}}, up to a 
subsequence, there exists $u\in H^s_r(\mathbb{R}^3)$ such that $u_n\rightharpoonup u$ 
weakly in $ H^s_r(\mathbb{R}^3)$, $u_n\to u$ strongly in $ L^p(\mathbb{R}^3)$, for 
$p\in (2, 2^*_s)$, $u_n(x)\rightharpoonup u(x)$ a.e.\ on $\mathbb{R}^3$. 
In view of $2+\frac{8s}{3}<p<2^*_s$, and Lemma {\ref{lem2.1311}} and  
Lemma {\ref{lem2.171}}, then
\begin{equation}\label{5.2122}
\lim_{n\to \infty}\int_{\mathbb{R}^3}|u_n|^pdx
=\int_{\mathbb{R}^3}|u|^pdx,~\lim_{n\to \infty}\int_{\mathbb{R}^3}
\phi^s_{u_n}u^2_ndx=\int_{\mathbb{R}^3}\phi^s_{u}u^2dx.
\end{equation}
Suppose by contradiction that, $u\equiv 0$. Then, by \eqref{5.2122} and 
$P_\mu(u_n)=o_n(1)$, we deduce that
\begin{align*}
o_n(1)&
= \Big(a+b\int_{\mathbb{R}^3}|(-\Delta)^{s/2}u_n|^2dx\Big)
 \int_{\mathbb{R}^3}|(-\Delta)^{s/2}u_n|^2dx
 +\frac{3-2s}{4s}\int_{\mathbb{R}^3}\phi^s_u u^2dx\\
&\quad -\mu \delta_{p,s}\int_{\mathbb{R}^3}|u_n|^pdx
 -\int_{\mathbb{R}^3}|u_n|^{2^*_s}dx\\
&= \Big(a+b\int_{\mathbb{R}^3}|(-\Delta)^{s/2}u_n|^2dx\Big)
 \int_{\mathbb{R}^3}|(-\Delta)^{s/2}u_n|^2dx
-\int_{\mathbb{R}^3}|u_n|^{2^*_s}dx+o_n(1).
\end{align*}
Without loss of generality, we assume that
$$
a\int_{\mathbb{R}^3}|(-\Delta)^{s/2}u_n|^2dx
+b\Big(\int_{\mathbb{R}^3}|(-\Delta)^{s/2}u_n|^2dx\Big)^2\to l\geq 0,\quad
\int_{\mathbb{R}^3}|u_n|^{2^*_s}dx\to l,
$$
as $n\to \infty$. If $l=0$, then we can deduce from the expression of $I_\mu(u_n)$ 
that $c_\mu(c)=0$, which is absurd since $c_\mu(c)>0$. Hence, $l>0$. 
By the definition of $S$, we have
$$
S\leq \frac{\int_{\mathbb{R}^3}|(-\Delta)^{s/2}u_n|^2dx}
{\|u_n\|^2_{2^*_s}}\leq \frac{1}{a}
\frac{a\int_{\mathbb{R}^3}|(-\Delta)^{s/2}u_n|^2dx
+b\Big(\int_{\mathbb{R}^3}|(-\Delta)^{s/2}u_n|^2dx\Big)^2}
{\left(\int_{\mathbb{R}^3}|u_n|^{2^*_s}
dx\right)^{2/2^*_s}}
\to \frac{1}{a}l^{\frac{2s}{3}}
$$
as $n\to \infty $. It follows that $l\geq (aS)^{\frac{3}{2s}}$. Consequently,
 by \eqref{5.2122} we have
\begin{align*}
c_\mu(c)
&= \lim_{n\to \infty}I_\mu(u_n)\\
&= \lim_{n\to \infty}\Big\{\frac{a}{2}\int_{\mathbb{R}^3}|(-\Delta)^{s/2}u_n|^2dx
 +\frac{b}{4}\Big(\int_{\mathbb{R}^3}|(-\Delta)^{s/2}u_n|^2dx\Big)^2 \\
&+ \frac{1}{4}\int_{\mathbb{R}^3}\phi^s_{u_n} u^2_ndx
  -\frac{\mu}{p}\int_{\mathbb{R}^3}|u_n|^pdx
  -\frac{1}{2^*_s}\int_{\mathbb{R}^3}|u_n|^{2^*_s}dx\Big\}\\
&\geq \Big(\frac{1}{4}-\frac{1}{2^*_s}\Big) l \\
&\geq \Big(\frac{1}{4}-\frac{1}{2^*_s}\Big)(aS)^{\frac{3}{2s}},
\end{align*}
which contradicts $I_\mu(u_n)\to c_\mu(c)
<\Big(\frac{1}{4}-\frac{1}{2^*_s}\Big)(aS)^{\frac{3}{2s}}$. 
Therefore, $u\not \equiv 0$.
\smallskip

\noindent\textbf{Step 2.} 
We prove that $u_n\to u$ in $L^{2^*_s}(\mathbb{R}^3)$.
Again by Lemma \ref{lem5.6}, we can obtain 
$\{\|(-\Delta)^{s/2}u_n\|_2\}$ is bounded in $\mathbb{R}$, by Prohorov's theorem 
\cite{P9899}, there exist two positive measures, $\zeta$,
 $\omega\in \mathcal{M}(\mathbb{R}^3)$, such that
\[ 
|(-\Delta)^{s/2}u_n|^2\rightharpoonup \omega, |u_n|^{2^*_s}\rightharpoonup\zeta \quad
\text{in } \mathcal{M}(\mathbb{R}^3)
\]
as $n\to \infty$. Then, by Lemma \ref{lem2.131}, either $u_n\to u$ in
$L^{2^*_s}_{loc}(\mathbb{R}^3)$ or there exists a (at most countable) set of distinct
points $\{x_j\}_{j\in J}\subset \mathbb{R}^3$ and positive numbers $\{\zeta_j\}_{j\in J}$ such that
$$
\zeta=|u|^{2^*_s}+\Sigma_{j\in J}\zeta_j \delta_{x_j}.
$$
Moreover, there exist some at most a countable set $J\subset \mathbb{N}$,
a corresponding set of distinct points $\{x_j\}_{j\in J}\subset \mathbb{R}^3$,
and two sets of positive numbers  $\{\zeta_j\}_{j\in J}\subset \mathbb{R}^3$ and
$\{\omega_j\}_{j\in J}\subset \mathbb{R}^3$ such that items
\eqref{lab3.11}-\eqref{lab3.13} hold. Now, assume that $J\neq \emptyset$.

 Similar to the proof in step 1 and step 2 of Lemma \ref{lem2.161},  we  
 can obtain
 \begin{equation}\label{4.1667}
 a\omega(\{x_j\})\leq\zeta_j,  \quad a\omega_\infty\leq\zeta_\infty.
 \end{equation}
Next, we claim that  $\zeta_j=0$ for any $j\in J$ and $\zeta_\infty=0$.

 Suppose by contradiction that, there exists $j_1\in J$ such that $\zeta_{j_1}>0$ 
or $\zeta_\infty>0$.  By  \eqref{4.1667}, Lemma {\ref{lem2.131}} and  
Lemma  \ref{lem2.11}, we obtain
$$
\zeta_{j_1}\geq (aS)^{\frac{3}{2s}} \quad\text{or}\quad
\zeta_{\infty}\geq (aS)^{\frac{3}{2s}}.
$$
If the former case occurs, one has
\begin{align*}
\Big(\frac{1}{4}-\frac{1}{2^*_s}\Big)(aS)^{\frac{3}{2s}}
&>c_\mu(c)
=\lim_{n\to \infty}\Big[I_\mu(u_n)-\frac{1}{4s}P_\mu(u_n)\Big]\\
&=  \lim_{n\to \infty}\big[\frac{1}{4}a\|(-\Delta)^{s/2}u_n\|^2_2
 +\Big(\frac{1}{4}-\frac{3-2s}{16s}\Big)\int_{\mathbb{R}^3}\phi^s_{u_n}u^2_ndx\\
&\quad +\Big(\frac{3(p-2)}{8ps}-\frac{1}{p}\Big)
 \mu \int_{\mathbb{R}^3}|u_n|^pdx
 +\Big(\frac{1}{4}-\frac{1}{2^*_s}\Big)\int_{\mathbb{R}^3}|u_n|^{2^*_s}dx\Big]\\
&=  \lim_{n\to \infty}\Big[\frac{1}{4}a\|(-\Delta)^{s/2}u_n\|^2_2
 +\frac{6s-3}{16s}\int_{\mathbb{R}^3}\phi^s_{u_n}u^2_ndx\\
&\quad +\frac{3(p-2)-8s}{8ps}\mu \int_{\mathbb{R}^3}|u_n|^pdx+\Big(\frac{1}{4}
 -\frac{1}{2^*_s}\Big)\int_{\mathbb{R}^3}|u_n|^{2^*_s}dx\Big]\\
&\geq \lim_{n\to \infty}\Big(\frac{1}{4}-\frac{1}{2^*_s}\Big)
 \int_{\mathbb{R}^3}|u_n|^{2^*_s}dx\\
&\geq \Big(\frac{1}{4}-\frac{1}{2^*_s}\Big)\lim_{\rho\to 0}
 \lim_{n\to \infty}\int_{\mathbb{R}^3}|u_n|^{2^*_s}\varphi_\rho dx\\
&=  \Big(\frac{1}{4}-\frac{1}{2^*_s}\Big)\zeta_{j_1}\\
&\geq\Big(\frac{1}{4}-\frac{1}{2^*_s}\Big) (aS)^{\frac{3}{2s}},
\end{align*}
which is a contradiction. If the latter case happens, we have
\begin{align*}
\Big(\frac{1}{4}-\frac{1}{2^*_s}\Big)(aS)^{\frac{3}{2s}}>c_\mu(c)
&\geq \lim_{n\to \infty}\Big(\frac{1}{4}-\frac{1}{2^*_s}\Big)
 \int_{\mathbb{R}^3}|u_n|^{2^*_s}dx\\
&\geq \Big(\frac{1}{4}-\frac{1}{2^*_s}\Big)\lim_{R\to \infty}
 \lim_{n\to \infty}\int_{\mathbb{R}^3}|u_n|^{2^*_s}\psi_R dx\\
&= \Big(\frac{1}{4}-\frac{1}{2^*_s}\Big)\zeta_\infty \\
&\geq\Big(\frac{1}{4}-\frac{1}{2^*_s}\Big) (aS)^{\frac{3}{2s}}.
\end{align*}
This is also a contradiction. Therefore, $\zeta_j=0$ for any 
$j\in J$ and $\zeta_\infty=0$. As a result, by Lemma \ref{lem2.131}, 
we obtain that $u_n\to u$ in $L^{2^*_s}_{loc}(\mathbb{R}^3)$.
Combining this with   Lemma \ref{lem2.11}, we know that $u_n\to u$ in 
$L^{2^*_s}(\mathbb{R}^3)$.
\smallskip

\noindent\textbf{Step 3.}
We prove that there exists some $\mu^*_2:=\mu^*_2(c)>0$ such that 
$\lim_{n\to +\infty}\lambda_n=\lambda<0$, if $\mu> \mu^*_2$ large.

By \eqref{5.16} and the fact that $u_n\in S_{r}(c)$, one has
\begin{align*}
&a\int_{\mathbb{R}^3}|(-\Delta)^{s/2}u_n|^2dx
 +b\Big(\int_{\mathbb{R}^3}|(-\Delta)^{s/2}u_n|^2dx\Big)^2
+\int_{\mathbb{R}^3}\phi^s_{u_n} u^2_ndx-\int_{\mathbb{R}^3}f(u_n)u_ndx\\
&= \lambda_n \int_{\mathbb{R}^3}|u_n|^2dx+o_n(1)\\
&= \lambda_n c^2+o_n(1).
\end{align*}
It indicates that
\begin{align*}
\lambda_n
&= \frac{1}{ c^2}\Big[a\int_{\mathbb{R}^3}|(-\Delta)^{s/2}u_n|^2dx
+b\Big(\int_{\mathbb{R}^3}|(-\Delta)^{s/2}u_n|^2dx\Big)^2\\
&\quad +\int_{\mathbb{R}^3}\phi^s_{u_n} u^2_ndx
 -\int_{\mathbb{R}^3}f(u_n)u_ndx\Big]+o_n(1).
\end{align*}
By Lemma \ref{lem5.6}, we  know that $\{u_n\}$ is bounded in $H^s_r(\mathbb{R}^3)$, 
and so, $\{\lambda_n\}$ is bounded in $\mathbb{R}$. 
By Lemma {\ref{lem5.8}} we know that 
$\limsup_{n\to \infty}|\lambda_n|\leq \frac{C}{c^2}c_\mu(c)$. 
From this and  $P_\mu(u_n)\to 0$ as $n\to \infty$, we derive that
\begin{align*}
\lambda_n
&= \frac{1}{ c^2}\Big[a\int_{\mathbb{R}^3}|(-\Delta)^{s/2}u_n|^2dx
 +b\Big(\int_{\mathbb{R}^3}|(-\Delta)^{s/2}u_n|^2dx\Big)^2\\
&\quad +\int_{\mathbb{R}^3}\phi^s_{u_n} u^2_ndx
 -\int_{\mathbb{R}^3}f(u_n)u_ndx-\frac{1}{s}P_\mu(u_n)\Big]+o_n(1)\\
&= \frac{1}{c^2}\Big[\frac{6s-3}{4s}\int_{\mathbb{R}^3}\phi^s_{u_n} u^2_ndx
 +\mu \frac{p(3-2s)-6}{2ps}\int_{\mathbb{R}^3}|u_n|^pdx\Big]+o_n(1).
\end{align*}
By \eqref{5.2122} and similar arguments to that of \eqref{4.32}-\eqref{4.35}, 
we see that there exists $\mu^*_2:=\mu^*_2(c)>0$, such that
\begin{equation}\label{5.212}
\begin{split}
\lambda&= \lim_{n\to \infty}\lambda_n\\
&= \lim_{n\to \infty} \frac{1}{c^2} \Big\{\frac{6s-3}{4s}
 \int_{\mathbb{R}^3}\phi^s_{u_n} u^2_ndx
 +\mu \frac{p(3-2s)-6}{2ps}\int_{\mathbb{R}^3}|u_n|^pdx\Big\}\\
&= \frac{1}{c^2} \Big\{\frac{6s-3}{4s}\int_{\mathbb{R}^3}\phi^s_{u} u^2dx
 +\mu \frac{p(3-2s)-6}{2ps}\int_{\mathbb{R}^3}|u|^pdx\Big\}
< 0,
\end{split}
\end{equation}
for $\mu>\mu^*_2$ large.

With the help  of  the above step 1, step 2, step 3, we can prove $u_n\to u$ 
strongly in $H^s_r(\mathbb{R}^3)$. Let $\mu>\mu^*_2$, set 
$B:=\lim_{n\to\infty} \|(-\Delta)^{s/2}u_n\|^2_2\geq 0 $, by the weak convergence of 
$u_n\rightharpoonup u$ in  $H^s_r(\mathbb{R}^3)$ and  \eqref{5.166}, one obtains
\begin{equation}\label{5.22}
\begin{split}
&(a+bB)\int_{\mathbb{R}^3}(-\Delta)^{s/2}u(-\Delta)^{s/2}\varphi dx
 +\int_{\mathbb{R}^3}\phi^s_{u}u\varphi dx\\
&-\mu \int_{\mathbb{R}^3}|u|^{q-2}u\varphi dx
 -\int_{\mathbb{R}^3}|u|^{2^*_s-2}u\varphi dx \\
&=\lambda \int_{\mathbb{R}^3} u\varphi dx.
\end{split}
\end{equation}
Therefore, from \eqref{5.2122}, \eqref{5.212}, and \eqref{5.22} and 
$u_n\to u$ in $L^{2^*_s}(\mathbb{R}^3)$, it follows that
\begin{align*} %\label{5.23}
&\lim_{n\to\infty}\Big[\Big(a+b\int_{\mathbb{R}^3}|(-\Delta)^{s/2}u_n|^2dx\Big)
 \int_{\mathbb{R}^3}|(-\Delta)^{s/2}u_n|^2 dx
 +\int_{\mathbb{R}^3}\phi^s_{u_n}{u^2_n}dx
 -\lambda \int_{\mathbb{R}^3} {u^2_n} dx\Big]\\
&=\lim_{n\to\infty}\Big [\mu \int_{\mathbb{R}^3}|u_n|^{p}dx
 +\int_{\mathbb{R}^3}|u_n|^{2^*_s} dx\Big]\\
&=\mu \int_{\mathbb{R}^3}|u|^{p}dx+\int_{\mathbb{R}^3}|u|^{2^*_s} dx\\
&=(a+bB)\int_{\mathbb{R}^3}|(-\Delta)^{s/2}u|^2 dx
 +\int_{\mathbb{R}^3}\phi^s_{u}{u^2}dx
 -\lambda \int_{\mathbb{R}^3} {u^2} dx.
\end{align*}
Since $\lambda<0$, as in the proof of Lemma {\ref{lem2.161}},  
we can derive that
\begin{align*} %\label{5.24}
\lim_{n\to\infty} \int_{\mathbb{R}^3} {u^2_n} dx
= \int_{\mathbb{R}^3} {u^2} dx,\quad
\lim_{n\to\infty}  \|(-\Delta)^{s/2}u_n\|^2_2 = \|(-\Delta)^{s/2}u\|^2_2.
\end{align*}
Therefore, $u_n\to u$ in $H^s_{r}(\mathbb{R}^3)$ and $\|u\|_2=c$. 
The proof is complete.
\end{proof}

With the help of the above technical lemmas, we can prove Theorem \ref{thm1.2} 
as follows.

